\question[2]
Identify the four parts of waves identified by the letters. 1/2 point per blank. 

\vspace*{2ex}

\begin{multicols}{2}

\textsc{a}.~\ifprintanswers\textbf{Wavelength}\else \rule{6cm}{0.4pt} \fi\\[1ex]

\textsc{b}.~\ifprintanswers\textbf{Wave height}\else \rule{6cm}{0.4pt} \fi\\[1ex]

\textsc{c}.~\ifprintanswers\textbf{Crest}\else \rule{6cm}{0.4pt} \fi\\[1ex]

\textsc{d}.~\ifprintanswers\textbf{Trough}\else \rule{6cm}{0.4pt} \fi

\columnbreak
	\includegraphics[width=0.47\textwidth]{wave_features_blank}

\end{multicols}

\question
	\begin{parts}
	\part[5]
	Illustrate the profile of a fringing reef from shoreline to open ocean. Label each zone.

	\AnswerBox{1\basespace}{%
	Fore-reef slope, reef crest, reef flat.%
}


	\part[10]
	Describe how light, aerial exposure, and wave action interact to determine the types of hermatypic corals (relative size, branching, etc.) present in each zone.
	
		\begin{AnswerPage}{5in}
		Answer goes here.
		\end{AnswerPage}
	\end{parts}


\question[10]
Describe how light, aerial exposure, and wave action interact to determine the types of hermatypic corals (relative size, branching, etc.) found on the reef crest.

\begin{AnswerPage}{0.45\textheight}
	Answer goes here.
\end{AnswerPage}



\end{questions}

